\documentclass[11pt]{article}
\usepackage[a4paper,margin=22mm]{geometry}
\usepackage{fontspec}
\setmainfont{LinLibertine_R.otf}[BoldFont=LinLibertine_RB.otf,ItalicFont=LinLibertine_RI.otf,BoldItalicFont=LinLibertine_RBI.otf]
\usepackage{xcolor}
\usepackage{hyperref}
\hypersetup{colorlinks=true,urlcolor=blue}
\setlength{\parindent}{0pt}
\setlength{\parskip}{6pt}
\sloppy
\emergencystretch=2em
\begin{document}
\section*{Historical Britannica: extracted source}
Encyclopaedia Britannica - Seventh edition, Volume 11, Grotius-HYD - EB.15\par
{\small Source row 2519. Parser: nvidia/nemotron-3.5-parse.\par}
{\small This is model-extracted text, not a corrected transcription. OCR errors, reading order and region alignment require human review.\par}
\hrule\medskip
EB 15
\par

{\scriptsize\color{gray} Region region-0; provider order 0.\par}
\medskip
\newpage\section*{Source record}
bri-2519-f242fdaba53d\par
Image SHA-256: \texttt{f242fdaba53d054518c22b5bea5686ee}\par\texttt{b887deeb1b2df944a177caf79d349c86}\par
Dataset revision: \texttt{309f765b1ecc83e0a7f7}\texttt{762e44b4a927a7555443}\par
National Library of Scotland, Encyclopaedia Britannica Illustrated (CC0).\par
\url{https://huggingface.co/datasets/NationalLibraryOfScotland/encyclopaedia_britannica_illustrated/viewer/default/train?row=2519}
\end{document}
